
\section*{CHƯƠNG 3 - THIẾT KẾ HỆ THỐNG}
\addcontentsline{toc}{section}{\numberline{}CHƯƠNG 3 - THIẾT KẾ HỆ THỐNG}
\setcounter{section}{3}
\setcounter{subsection}{0}
\setcounter{figure}{0}
\setcounter{table}{0}

\subsection{Tổng quan kiến trúc hệ thống}

\subsubsection{Kiến trúc tổng thể}
Hệ thống phân tích cảm xúc được thiết kế theo mô hình pipeline với 6 bước xử lý tuần tự. Theo yêu cầu đề bài, dự án tập trung vào 3 bước chính:

\begin{enumerate}
    \item \textbf{Thu thập dữ liệu:} Sử dụng YouTube Data API v3 (lướt qua)
    \item \textbf{Chuẩn hóa dữ liệu:} Làm sạch và tiền xử lý văn bản (lướt qua)
    \item \textbf{Gán nhãn:} Phân tích cảm xúc với TextBlob (lướt qua)
    \item \textbf{Thực nghiệm mô hình:} Phân tích mối quan hệ các đặc trưng
    \item \textbf{Xây dựng ứng dụng:} Demo phân tích văn bản mới
    \item \textbf{Phân tích sâu kết quả:} Thống kê và insights
\end{enumerate}

\subsubsection{Công nghệ sử dụng}
\begin{itemize}
    \item \textbf{Ngôn ngữ:} Python 3.7+
    \item \textbf{API:} YouTube Data API v3
    \item \textbf{NLP:} TextBlob
    \item \textbf{Data Processing:} Pandas, NumPy
    \item \textbf{Visualization:} Matplotlib, Seaborn
\end{itemize}

\subsection{Thiết kế module thu thập và chuẩn hóa (Tóm tắt)}

\subsubsection{Thu thập dữ liệu}
Hệ thống sử dụng YouTube Data API v3 với 4 API keys dự phòng để thu thập 100,000 bình luận từ video "Avengers: Endgame". Dữ liệu bao gồm: tên phim, năm phát hành, nội dung bình luận, tác giả, số likes, và thời gian đăng.

\subsubsection{Chuẩn hóa dữ liệu}
Quy trình làm sạch 8 bước: loại bỏ URL, mentions, hashtags, email, số điện thoại, ký tự đặc biệt, chuẩn hóa khoảng trắng, và chuyển chữ thường. Sau đó lọc bỏ bình luận rỗng, quá ngắn, và trùng lặp. Kết quả: 92,521 bình luận hợp lệ từ 100,000 ban đầu.

\subsubsection{Gán nhãn}
Sử dụng TextBlob để tính điểm polarity (-1 đến 1) và phân loại thành Positive/Negative/Neutral với 3 mức độ tin cậy (High/Medium/Low). Mỗi bình luận có 5 thuộc tính: sentiment\_score, subjectivity, sentiment\_label, confidence\_level, và sentiment\_description.

\subsection{Thiết kế module thực nghiệm mô hình (TẬP TRUNG)}

\subsubsection{Mục tiêu thực nghiệm}
Phân tích mối quan hệ giữa các đặc trưng để trả lời các câu hỏi nghiên cứu:

\begin{enumerate}
    \item Độ dài bình luận có ảnh hưởng đến cảm xúc không?
    \item Bình luận nào nhận được nhiều likes hơn?
    \item Phân bố điểm cảm xúc như thế nào?
    \item Các đặc trưng có tương quan với nhau không?
\end{enumerate}

\subsubsection{Phương pháp thực nghiệm}
\textbf{1. Phân tích tương quan (Correlation Analysis):}
\begin{itemize}
    \item Tính ma trận tương quan Pearson giữa sentiment\_score, cleaned\_length, likes
    \item Trực quan hóa bằng heatmap
    \item Giải thích ý nghĩa các hệ số tương quan
\end{itemize}

\textbf{2. Phân tích phân bố (Distribution Analysis):}
\begin{itemize}
    \item Histogram: Phân bố điểm cảm xúc
    \item Box Plot: Phân bố điểm theo từng nhãn (Positive/Negative/Neutral)
    \item Phát hiện outliers và giải thích
\end{itemize}

\textbf{3. Phân tích so sánh (Comparative Analysis):}
\begin{itemize}
    \item So sánh độ dài trung bình giữa các nhãn
    \item So sánh số likes trung bình giữa các nhãn
    \item Kiểm định thống kê (t-test, ANOVA)
\end{itemize}

\textbf{4. Phân tích hồi quy (Regression Analysis):}
\begin{itemize}
    \item Scatter plot: Độ dài vs Cảm xúc
    \item Scatter plot: Likes vs Cảm xúc
    \item Tìm xu hướng và pattern
\end{itemize}

\subsubsection{Các biểu đồ thực nghiệm}
\begin{table}[H]
\centering
\caption{Danh sách biểu đồ thực nghiệm}
\begin{tabular}{|l|l|p{6cm}|}
\hline
\textbf{STT} & \textbf{Loại biểu đồ} & \textbf{Mục đích} \\
\hline
1 & Bar Chart & Phân bố số lượng theo nhãn \\
\hline
2 & Pie Chart & Tỷ lệ phần trăm các nhãn \\
\hline
3 & Histogram & Phân bố điểm cảm xúc liên tục \\
\hline
4 & Scatter Plot (Length vs Sentiment) & Mối quan hệ độ dài - cảm xúc \\
\hline
5 & Scatter Plot (Likes vs Sentiment) & Mối quan hệ likes - cảm xúc \\
\hline
6 & Box Plot & Phân bố điểm theo nhãn, phát hiện outliers \\
\hline
7 & Horizontal Bar & Độ dài trung bình theo nhãn \\
\hline
8 & Heatmap & Ma trận tương quan giữa các đặc trưng \\
\hline
\end{tabular}
\end{table}

\subsubsection{Metrics đánh giá}
\begin{itemize}
    \item \textbf{Pearson Correlation Coefficient (r):} Đo lường tương quan tuyến tính (-1 đến 1)
    \item \textbf{Mean, Median, Std:} Thống kê mô tả
    \item \textbf{Quartiles (Q1, Q2, Q3):} Phân vị để phát hiện outliers
    \item \textbf{IQR (Interquartile Range):} Độ phân tán
\end{itemize}

\subsection{Thiết kế cấu trúc dữ liệu}

\subsubsection{File youtube\_raw\_data.csv}
Dữ liệu thô sau khi thu thập (6 cột):

\begin{verbatim}
    movie_name, release_year, review_text, author, 
    likes, published_at
\end{verbatim}

\subsubsection{File youtube\_clean\_data.csv}
Dữ liệu sau khi chuẩn hóa (10 cột):

\begin{verbatim}
    movie_name, release_year, review_text, author, 
    likes, published_at, cleaned_text, original_length, 
    cleaned_length, reduction_percent
\end{verbatim}

\subsubsection{File youtube\_labeled\_data.csv}
Dữ liệu sau khi gán nhãn (12 cột):

\begin{verbatim}
    movie_name, release_year, review_text, cleaned_text,
    sentiment_score, subjectivity, sentiment_label,
    confidence_level, sentiment_description, likes,
    author, published_at
\end{verbatim}

\subsection{Thiết kế ứng dụng demo}

\subsubsection{Tổng quan ứng dụng}
Ứng dụng demo là một công cụ phân tích cảm xúc văn bản thời gian thực, cho phép người dùng nhập bất kỳ văn bản tiếng Anh nào và nhận được kết quả phân tích ngay lập tức. Ứng dụng được thiết kế với giao diện console đơn giản, dễ sử dụng, phù hợp cho mục đích học tập và demo.

\textbf{Mục tiêu ứng dụng:}
\begin{itemize}
    \item Cung cấp công cụ phân tích cảm xúc nhanh chóng và chính xác
    \item Giúp người dùng hiểu rõ quá trình xử lý văn bản
    \item Demo các kỹ thuật NLP và sentiment analysis
    \item Hỗ trợ nghiên cứu và thử nghiệm với dữ liệu mới
\end{itemize}

\subsubsection{Kiến trúc ứng dụng}
Ứng dụng được thiết kế theo mô hình 3 lớp:

\textbf{1. Input Layer (Lớp nhập liệu):}
\begin{itemize}
    \item Nhận văn bản tiếng Anh từ người dùng qua console
    \item Validate input: kiểm tra không rỗng, độ dài hợp lệ (> 3 ký tự)
    \item Hỗ trợ 2 chế độ: Single text (phân tích 1 câu) và Batch processing (nhiều câu)
    \item Cho phép người dùng thoát bằng lệnh 'quit' hoặc 'exit'
\end{itemize}

\textbf{2. Processing Layer (Lớp xử lý):}
\begin{itemize}
    \item \textbf{Text Cleaning:} Áp dụng 8 bước làm sạch (loại bỏ URL, mentions, hashtags, email, số điện thoại, ký tự đặc biệt, chuẩn hóa khoảng trắng, chuyển chữ thường)
    \item \textbf{Sentiment Analysis:} Tính polarity score (-1 đến 1) và subjectivity (0 đến 1) bằng TextBlob
    \item \textbf{Classification:} Phân loại thành Positive (polarity > 0.1), Negative (polarity < -0.1), hoặc Neutral
    \item \textbf{Confidence Scoring:} Xác định độ tin cậy High (|polarity| > 0.5), Medium (0.1 < |polarity| ≤ 0.5), Low (|polarity| ≤ 0.1)
\end{itemize}

\textbf{3. Output Layer (Lớp hiển thị):}
\begin{itemize}
    \item Hiển thị kết quả với emoji và màu sắc (😊 Positive, 😐 Neutral, 😢 Negative)
    \item Thống kê chi tiết quá trình làm sạch (số ký tự trước/sau, tỷ lệ giảm)
    \item Điểm số sentiment và độ tin cậy
    \item So sánh văn bản gốc và văn bản đã làm sạch
    \item Mô tả cảm xúc bằng ngôn ngữ tự nhiên
\end{itemize}

\subsubsection{Các chức năng chính}

\textbf{Chức năng 1: Phân tích văn bản đơn}
\begin{itemize}
    \item \textbf{Input:} 1 câu văn bản tiếng Anh
    \item \textbf{Process:} Làm sạch → Phân tích → Phân loại
    \item \textbf{Output:} Nhãn cảm xúc, điểm số, độ tin cậy, thống kê chi tiết
    \item \textbf{Use case:} Test nhanh với câu ngắn, demo cho người dùng mới
    \item \textbf{Thời gian xử lý:} < 0.5 giây
\end{itemize}

\textbf{Chức năng 2: Phân tích batch}
\begin{itemize}
    \item \textbf{Input:} Danh sách nhiều văn bản (từ file hoặc nhập liên tục)
    \item \textbf{Process:} Xử lý tuần tự từng văn bản
    \item \textbf{Output:} Bảng tổng hợp kết quả với các cột: Text, Label, Score, Confidence
    \item \textbf{Use case:} Phân tích nhiều bình luận cùng lúc, so sánh kết quả
    \item \textbf{Thời gian xử lý:} ~0.3 giây/văn bản
\end{itemize}

\textbf{Chức năng 3: So sánh văn bản}
\begin{itemize}
    \item \textbf{Input:} 2 văn bản để so sánh
    \item \textbf{Process:} Phân tích cả 2 văn bản độc lập
    \item \textbf{Output:} So sánh điểm số, nhãn, độ dài, độ tin cậy
    \item \textbf{Use case:} So sánh cảm xúc giữa 2 bình luận, phân tích sự khác biệt
    \item \textbf{Highlight:} Hiển thị văn bản nào tích cực/tiêu cực hơn
\end{itemize}

\textbf{Chức năng 4: Thống kê chi tiết}
\begin{itemize}
    \item Hiển thị quá trình làm sạch từng bước với số lượng ký tự loại bỏ
    \item Số URLs, mentions, hashtags, emails, số điện thoại đã loại bỏ
    \item Tỷ lệ giảm độ dài (reduction percentage)
    \item Phân tích độ chủ quan (subjectivity) của văn bản
\end{itemize}

\textbf{Chức năng 5: Lưu kết quả}
\begin{itemize}
    \item Lưu kết quả phân tích vào file CSV
    \item Các cột: timestamp, original\_text, cleaned\_text, sentiment\_label, sentiment\_score, confidence\_level
    \item Cho phép người dùng xem lại lịch sử phân tích
    \item Hỗ trợ export để phân tích thêm bằng Excel hoặc Python
\end{itemize}

\subsubsection{Giao diện người dùng}

Ứng dụng demo được chia thành 2 phần chính để phục vụ các mục đích khác nhau:

\textbf{PHẦN 1: Test với mẫu có sẵn}

Phần này tự động chạy 5 mẫu văn bản được định sẵn để demo các trường hợp điển hình:

\begin{enumerate}
    \item Bình luận rất tích cực: "This movie is absolutely amazing! Best film ever!"
    \item Bình luận tiêu cực: "Terrible movie, waste of time."
    \item Bình luận trung lập: "It was okay, nothing special."
    \item Bình luận tích cực: "I love this so much! Can't wait to watch it again!"
    \item Bình luận rất tiêu cực: "Worst experience ever. Very disappointed."
\end{enumerate}

\textbf{PHẦN 2: Chế độ tương tác (Interactive Mode)}

Sau khi xem 5 mẫu, người dùng có thể tự nhập văn bản để phân tích:

\begin{itemize}
    \item \textbf{Tính năng:} Cho phép người dùng nhập bất kỳ văn bản tiếng Anh nào
    \item \textbf{Số lượng:} Không giới hạn số lần phân tích
    \item \textbf{Thoát:} Nhập 'quit', 'exit', hoặc 'q' để kết thúc
    \item \textbf{Tiếp tục:} Sau mỗi lần phân tích, hỏi người dùng có muốn tiếp tục không
    \item \textbf{Validation:} Kiểm tra độ dài tối thiểu (3 ký tự)
    \item \textbf{Error handling:} Xử lý các trường hợp lỗi (văn bản rỗng, quá ngắn)
\end{itemize}

\textbf{Console Interface (CLI):}

\textit{Phần 1 - Mẫu có sẵn:}
\begin{verbatim}
    ======================================================================
    🎬 BƯỚC 6: ỨNG DỤNG DEMO
    ======================================================================
    
    🧪 PHẦN 1: Test với 5 mẫu văn bản có sẵn:
    
    ──────────────────────────────────────────────────────────────────────
    Mẫu 1/5:
    
    ======================================================================
    📝 Văn bản gốc:
       This movie is absolutely amazing! Best film ever!
    
    🧹 Văn bản đã làm sạch:
       this movie is absolutely amazing best film ever
    
    📊 Thống kê làm sạch:
       • Độ dài: 49 → 46 ký tự
       • Giảm: 6.1%
       • URLs loại bỏ: 0
       • Mentions loại bỏ: 0
    
    🎯 Kết quả:
       😊 Nhãn: Positive
       📊 Mô tả: Rất tích cực
       🎚️  Độ tin cậy: High
       📈 Điểm cảm xúc: 0.850
       📊 Độ chủ quan: 0.900
    ======================================================================
    
    [... 4 mẫu còn lại ...]
\end{verbatim}

\textit{Phần 2 - Tự nhập:}
\begin{verbatim}
    ======================================================================
    🎯 PHẦN 2: TỰ NHẬP VĂN BẢN ĐỂ PHÂN TÍCH
    ======================================================================
    
    💡 Bạn có thể nhập bất kỳ văn bản tiếng Anh nào để phân tích cảm xúc.
       Nhập 'quit', 'exit', hoặc 'q' để kết thúc.
    
    ──────────────────────────────────────────────────────────────────────
    📝 Nhập văn bản của bạn (hoặc 'quit' để thoát): I love this movie!
    
    🔍 Đang phân tích văn bản #1...
    
    ======================================================================
    📝 Văn bản gốc:
       I love this movie!
    
    🧹 Văn bản đã làm sạch:
       i love this movie
    
    📊 Thống kê làm sạch:
       • Độ dài: 18 → 17 ký tự
       • Giảm: 5.6%
       • URLs loại bỏ: 0
       • Mentions loại bỏ: 0
    
    🎯 Kết quả:
       😊 Nhãn: Positive
       📊 Mô tả: Tích cực
       🎚️  Độ tin cậy: Medium
       📈 Điểm cảm xúc: 0.500
       📊 Độ chủ quan: 0.600
    ======================================================================
    
    ──────────────────────────────────────────────────────────────────────
    ❓ Bạn có muốn phân tích văn bản khác không? (y/n): y
    
    ──────────────────────────────────────────────────────────────────────
    📝 Nhập văn bản của bạn (hoặc 'quit' để thoát): This is terrible!
    
    🔍 Đang phân tích văn bản #2...
    
    ======================================================================
    📝 Văn bản gốc:
       This is terrible!
    
    🎯 Kết quả:
       😞 Nhãn: Negative
       📊 Mô tả: Tiêu cực
       🎚️  Độ tin cậy: High
       📈 Điểm cảm xúc: -0.700
       📊 Độ chủ quan: 0.800
    ======================================================================
    
    ──────────────────────────────────────────────────────────────────────
    ❓ Bạn có muốn phân tích văn bản khác không? (y/n): n
    
    ✅ Đã phân tích tổng cộng 2 văn bản của bạn.
    
    ======================================================================
    🎉 HOÀN TẤT DEMO ỨNG DỤNG!
    ======================================================================
\end{verbatim}

\textbf{Ưu điểm của chế độ tương tác:}
\begin{itemize}
    \item \textbf{Linh hoạt:} Người dùng có thể test với bất kỳ văn bản nào
    \item \textbf{Thực tế:} Mô phỏng cách sử dụng thực tế của ứng dụng
    \item \textbf{Học tập:} Giúp người dùng hiểu rõ cách hoạt động của thuật toán
    \item \textbf{So sánh:} Có thể nhập nhiều văn bản để so sánh kết quả
    \item \textbf{Không giới hạn:} Phân tích bao nhiêu văn bản tùy thích
\end{itemize}

\textbf{Các trường hợp xử lý đặc biệt:}
\begin{itemize}
    \item \textbf{Văn bản rỗng:} Hiển thị "⚠️ Bạn chưa nhập văn bản nào" và yêu cầu nhập lại
    \item \textbf{Văn bản quá ngắn:} Cảnh báo nhưng vẫn phân tích với độ tin cậy Low
    \item \textbf{Ngắt giữa chừng:} Hỗ trợ Ctrl+C để thoát bất kỳ lúc nào
    \item \textbf{Đếm số lượng:} Hiển thị số văn bản đã phân tích khi kết thúc
\end{itemize}

\subsubsection{Xử lý edge cases}

\textbf{Case 1: Văn bản rỗng}
\begin{itemize}
    \item Kiểm tra và báo lỗi: "⚠️ Vui lòng nhập văn bản"
    \item Yêu cầu nhập lại
    \item Không thực hiện phân tích
\end{itemize}

\textbf{Case 2: Văn bản quá ngắn (< 3 ký tự)}
\begin{itemize}
    \item Cảnh báo: "⚠️ Văn bản quá ngắn, kết quả có thể không chính xác"
    \item Vẫn phân tích nhưng đánh dấu Low confidence
    \item Hiển thị gợi ý: "Nên nhập ít nhất 10 ký tự để có kết quả tốt hơn"
\end{itemize}

\textbf{Case 3: Văn bản không phải tiếng Anh}
\begin{itemize}
    \item Phát hiện ngôn ngữ bằng langdetect (nếu có)
    \item Cảnh báo: "⚠️ TextBlob hoạt động tốt nhất với tiếng Anh"
    \item Vẫn thực hiện phân tích nhưng kết quả có thể không chính xác
\end{itemize}

\textbf{Case 4: Văn bản chỉ có emoji/số}
\begin{itemize}
    \item Sau làm sạch trở thành rỗng hoặc quá ngắn
    \item Trả về Neutral với Low confidence
    \item Thông báo: "Văn bản không chứa nội dung văn bản có nghĩa"
\end{itemize}

\textbf{Case 5: Văn bản quá dài (> 1000 ký tự)}
\begin{itemize}
    \item Cảnh báo: "Văn bản dài, có thể mất vài giây để xử lý"
    \item Vẫn xử lý bình thường
    \item Có thể cắt ngắn để hiển thị trong output
\end{itemize}

\subsubsection{Ứng dụng thực tế}

\textbf{1. Phân tích phản hồi khách hàng:}
\begin{itemize}
    \item Doanh nghiệp có thể sử dụng để phân tích feedback từ khách hàng
    \item Tự động phân loại đánh giá tích cực/tiêu cực
    \item Ưu tiên xử lý các phản hồi tiêu cực
    \item Thống kê tỷ lệ hài lòng của khách hàng
\end{itemize}

\textbf{2. Giám sát mạng xã hội:}
\begin{itemize}
    \item Theo dõi cảm xúc của người dùng về sản phẩm/dịch vụ
    \item Phát hiện sớm các vấn đề tiêu cực
    \item Đo lường hiệu quả chiến dịch marketing
    \item Phân tích xu hướng cảm xúc theo thời gian
\end{itemize}

\textbf{3. Nghiên cứu thị trường:}
\begin{itemize}
    \item Phân tích ý kiến người dùng về đối thủ cạnh tranh
    \item Tìm hiểu điểm mạnh/yếu của sản phẩm
    \item Xác định nhu cầu và mong muốn của khách hàng
    \item Hỗ trợ ra quyết định kinh doanh
\end{itemize}

\textbf{4. Hỗ trợ dịch vụ khách hàng:}
\begin{itemize}
    \item Tự động phân loại và định tuyến yêu cầu hỗ trợ
    \item Ưu tiên xử lý các khiếu nại nghiêm trọng
    \item Đánh giá chất lượng dịch vụ qua phản hồi
    \item Cải thiện quy trình chăm sóc khách hàng
\end{itemize}

\textbf{5. Phân tích nội dung giải trí:}
\begin{itemize}
    \item Đánh giá phản ứng của khán giả với phim/nhạc/game
    \item Dự đoán thành công của sản phẩm giải trí
    \item Phân tích xu hướng và sở thích người dùng
    \item Hỗ trợ sản xuất nội dung phù hợp
\end{itemize}

\subsubsection{Ví dụ sử dụng thực tế}

\textbf{Ví dụ 1: Phân tích bình luận tích cực}

\textit{Input:}
\begin{verbatim}
    "This is the best movie I've ever seen! 
     Amazing acting and great story!"
\end{verbatim}

\textit{Output:}
\begin{itemize}
    \item Nhãn: 😊 Positive
    \item Điểm cảm xúc: 0.85
    \item Độ tin cậy: High
    \item Mô tả: Rất tích cực
    \item Độ chủ quan: 0.90 (rất chủ quan)
\end{itemize}

\textbf{Ví dụ 2: Phân tích bình luận tiêu cực}

\textit{Input:}
\begin{verbatim}
    "Terrible movie, waste of time and money. 
     Very disappointed."
\end{verbatim}

\textit{Output:}
\begin{itemize}
    \item Nhãn: 😢 Negative
    \item Điểm cảm xúc: -0.72
    \item Độ tin cậy: High
    \item Mô tả: Rất tiêu cực
    \item Độ chủ quan: 0.85
\end{itemize}

\textbf{Ví dụ 3: Phân tích bình luận trung lập}

\textit{Input:}
\begin{verbatim}
    "The movie was released in 2019. 
     It has a runtime of 181 minutes."
\end{verbatim}

\textit{Output:}
\begin{itemize}
    \item Nhãn: 😐 Neutral
    \item Điểm cảm xúc: 0.02
    \item Độ tin cậy: Low
    \item Mô tả: Trung lập
    \item Độ chủ quan: 0.15 (khách quan)
\end{itemize}

\subsubsection{Đánh giá ứng dụng}

\textbf{Ưu điểm:}
\begin{itemize}
    \item \textbf{Đơn giản, dễ sử dụng:} Giao diện console trực quan, không cần cài đặt phức tạp
    \item \textbf{Phản hồi nhanh:} Xử lý và trả kết quả trong < 1 giây
    \item \textbf{Hiển thị chi tiết:} Cung cấp đầy đủ thông tin về quá trình xử lý
    \item \textbf{Hỗ trợ nhiều chế độ:} Single, batch, compare, statistics
    \item \textbf{Xử lý edge cases tốt:} Kiểm tra và cảnh báo các trường hợp đặc biệt
    \item \textbf{Lưu kết quả:} Hỗ trợ export sang CSV để phân tích thêm
    \item \textbf{Miễn phí và mã nguồn mở:} Dễ dàng tùy chỉnh và mở rộng
\end{itemize}

\textbf{Hạn chế:}
\begin{itemize}
    \item \textbf{Chỉ hỗ trợ CLI:} Chưa có giao diện đồ họa (GUI), khó sử dụng với người dùng không quen command line
    \item \textbf{Chưa lưu lịch sử:} Không có database để lưu trữ và tra cứu kết quả cũ
    \item \textbf{Chưa có API endpoint:} Không thể tích hợp vào hệ thống khác qua REST API
    \item \textbf{Chỉ hỗ trợ tiếng Anh:} TextBlob hoạt động tốt nhất với tiếng Anh, các ngôn ngữ khác kém chính xác
    \item \textbf{Phụ thuộc TextBlob:} Độ chính xác phụ thuộc vào thuật toán của TextBlob
    \item \textbf{Không có authentication:} Chưa có hệ thống quản lý người dùng
\end{itemize}

\textbf{So sánh với các công cụ khác:}

\begin{table}[H]
\centering
\caption{So sánh ứng dụng demo với các công cụ khác}
\begin{tabular}{|l|c|c|c|}
\hline
\textbf{Tiêu chí} & \textbf{Demo này} & \textbf{VADER} & \textbf{Google NLP API} \\
\hline
Dễ sử dụng & ⭐⭐⭐⭐⭐ & ⭐⭐⭐⭐ & ⭐⭐⭐ \\
\hline
Tốc độ & ⭐⭐⭐⭐⭐ & ⭐⭐⭐⭐⭐ & ⭐⭐⭐ \\
\hline
Độ chính xác & ⭐⭐⭐ & ⭐⭐⭐⭐ & ⭐⭐⭐⭐⭐ \\
\hline
Chi phí & Miễn phí & Miễn phí & Trả phí \\
\hline
Tùy chỉnh & ⭐⭐⭐⭐⭐ & ⭐⭐⭐⭐ & ⭐⭐ \\
\hline
Hỗ trợ ngôn ngữ & Tiếng Anh & Tiếng Anh & 100+ ngôn ngữ \\
\hline
\end{tabular}
\end{table}

\subsubsection{Hướng phát triển}

\textbf{Ngắn hạn (1-3 tháng):}
\begin{itemize}
    \item Xây dựng giao diện web đơn giản với Flask/Django
    \item Thêm database SQLite để lưu lịch sử phân tích
    \item Hỗ trợ upload file CSV/TXT để phân tích batch
    \item Thêm biểu đồ trực quan hóa kết quả (pie chart, bar chart)
    \item Cải thiện xử lý edge cases và validation
\end{itemize}

\textbf{Trung hạn (3-6 tháng):}
\begin{itemize}
    \item Xây dựng REST API với FastAPI
    \item Thêm authentication và authorization
    \item Hỗ trợ nhiều mô hình sentiment analysis (VADER, BERT, RoBERTa)
    \item Cho phép người dùng chọn mô hình phù hợp
    \item Thêm tính năng so sánh hiệu năng các mô hình
    \item Deploy lên cloud (Heroku, AWS, Google Cloud)
\end{itemize}

\textbf{Dài hạn (6-12 tháng):}
\begin{itemize}
    \item Xây dựng mobile app (React Native, Flutter)
    \item Hỗ trợ phân tích real-time từ social media (Twitter, Facebook)
    \item Thêm tính năng phân tích aspect-based sentiment
    \item Tích hợp machine learning để cải thiện độ chính xác
    \item Hỗ trợ đa ngôn ngữ (tiếng Việt, tiếng Trung, tiếng Nhật)
    \item Xây dựng dashboard analytics với insights tự động
\end{itemize}

\subsection{Tổng kết}
Chương này đã trình bày thiết kế chi tiết của hệ thống, bao gồm:

\begin{itemize}
    \item Kiến trúc tổng thể với 6 bước xử lý
    \item Module thu thập dữ liệu với hệ thống API dự phòng
    \item Module chuẩn hóa với 8 bước làm sạch
    \item Module gán nhãn với 3 mức độ tin cậy
    \item Module phân tích và trực quan hóa
    \item Cấu trúc dữ liệu đầu vào/đầu ra
    \item Thiết kế ứng dụng demo
\end{itemize}

Thiết kế này đảm bảo hệ thống hoạt động hiệu quả, dễ bảo trì và mở rộng. Các phương pháp và công nghệ được lựa chọn dựa trên nghiên cứu của \cite{hutto2014vader}, \cite{liu2012sentiment}, và \cite{pang2008opinion}.
